% Generated draft: requires mathematical and source review.
\documentclass[11pt]{article}
\usepackage[margin=25mm]{geometry}
\usepackage{fontspec}
\setmainfont{Latin Modern Roman}
\usepackage{amsmath,amssymb,mathtools,unicode-math}
\setmathfont{Latin Modern Math}
\usepackage{enumitem,needspace,xurl,hyperref}
\setlength{\parindent}{0pt}
\setlength{\parskip}{5pt}
\setlength{\emergencystretch}{3em}
\newcommand{\N}{\mathbb{N}}
\newcommand{\Z}{\mathbb{Z}}
\newcommand{\Q}{\mathbb{Q}}
\newcommand{\R}{\mathbb{R}}
\newcommand{\C}{\mathbb{C}}
\newcommand{\F}{\mathbb{F}}
\newcommand{\A}{\mathbb{A}}
\newcommand{\PP}{\mathbb{P}}
\newcommand{\E}{\mathbb{E}}
\newcommand{\eps}{\varepsilon}
\newcommand{\dd}{\,\mathrm{d}}
\begin{document}
\section*{UnsolvedMath ID: 30006563}
\textbf{Source problem code:} catalogue entry
\subsection{Definitions and mathematical statement}
For each positive integer \(n\), let \(K_n\) be the complete simple graph on vertex set \([n]=\{1,\ldots,n\}\). An edge-coloring with at most \(q\) colors is a map \(\chi:E(K_n)\to[q]\), where \(q\) is a nonnegative integer and \([0]=\varnothing\). No properness condition is imposed. Two triangles on disjoint three-element vertex sets \(A,B\subseteq[n]\) are \emph{color-isomorphic} if there is a bijection \(f:A\to B\) such that
\[
\chi(\{u,v\})=\chi(\{f(u),f(v)\})
\qquad\text{for every }\{u,v\}\in\binom{A}{2}.
\]
Thus the bijection preserves the actual color labels. Define
\[
g(n)=\min\bigl\{q\in\Z_{\geq0}:\text{there exists an edge-coloring }\chi:E(K_n)\to[q]
\text{ with no two vertex-disjoint color-isomorphic triangles}\bigr\}.
\]
The convention allowing \(q=0\) handles the edgeless graph \(K_1\). For \(n<6\), the forbidden configuration cannot occur. The source does not explicitly specify these small-order conventions or define color-isomorphism; the definitions above make the usual color-preserving interpretation explicit.

\textbf{Problem (David Conlon, Question 14).} Determine the asymptotic order of \(g(n)\) as \(n\to\infty\). In particular, is
\[
g(n)=\Theta(n),
\]
meaning that there exist constants \(a,b>0\) and an integer \(N\geq1\) such that \(an\leq g(n)\leq bn\) for every integer \(n\geq N\)? More strongly, is
\[
g(n)=n-O(1),
\]
meaning that there exist a constant \(C\geq0\) and an integer \(N\geq1\) such that \(|g(n)-n|\leq C\) for every integer \(n\geq N\)?

\textit{Supplied context and qualifications.} The supplied dated assessment describes bounds between order \(\sqrt n\) and \(n-4\). The original excerpt contains a displaced square-root symbol and a corrupted lower-bound display, so that display is not reliable as transcribed. The excerpt additionally attributes to Conlon, Fox and Pham an improvement of the form
\[
g(n)\geq c n^{1/2+\varepsilon}
\]
for some constants \(c>0\) and \(\varepsilon>0\). These bounds are recorded as supplied contextual claims, with no independently verified range of validity; in particular, the displayed upper bound \(n-4\) cannot apply to all positive integers \(n\). The source labels the problem open and describes both linear growth and the stronger bounded-distance assertion as unproved. These are supplied status assessments, not independent verification.
\subsection{Short English statement}
Color the edges of a complete graph so that no two triangles with disjoint vertex sets have exactly matching edge colors under a bijection of their vertices. How many colors are necessary as the number of vertices grows? The supplied question asks whether this minimum grows linearly, and even whether it stays within a constant of the number of vertices. The source reports the problem as open; its contextual bounds and status require review.
\subsection{Sources}
\begin{itemize}
\item Local UnsolvedMath registry, supplied complete source record: permanent ID \(30006563\), source problem code \texttt{catalogue entry}, title \emph{Colorings Without Disjoint Color-Isomorphic Triangles}; collection: Oberwolfach Reports Open Problems. The original excerpt attributes the question to David Conlon as Question 14, although the structured proposer field is empty.
\item Supplied report citation: \emph{Combinatorics} (2026), Oberwolfach Reports, DOI 10.4171/owr/2026/1, pp. 5--79. \url{https://doi.org/10.4171/owr/2026/1}. The supplied background describes Oberwolfach Reports as documenting mathematical workshops and their open-problem discussions.
\item Supplied literature source: \emph{Oberwolfach Report 1/2026}, \url{https://doi.org/10.4171/owr/2026/1}. Literature checked at: \texttt{2026-08-22}. This is supplied dated curation, not independently verified. Its assessment reports square-root-order to \(n-4\) bounds and states that linear growth and \(n-O(1)\) behavior remain unproved; the original excerpt also attributes a stronger lower bound to Conlon, Fox and Pham without a separate bibliographic reference.
\item The registry labels the statement \texttt{corrected\_verified} and says that the official report was checked, the question isolated from surrounding exposition, and its notation repaired. These are supplied verification claims; no source was visited or independently checked in preparing this draft.
\end{itemize}
\end{document}
